% !TeX root = main.tex
\chapter{はじまり}
\textbf{心の満足について}

誰にも妨げられないものが自由である。(87)

一人でいるときには、自分は神的なものに似ていると思い、多くの人といるときはお祭りとよび、そのようにしてすべてを満足する。（P90）

自分の意志に反しているところ、それが牢獄なのだ。（90）

全体に比べれば、我々はとても小さな部分である。しかしこれは肉体に関してのことである。理性に関しては神々になんら劣りもしなければ、より小さくもない。(91)

神々は何に対して君に責任があるようにしたのか、それは、君の力の及ぶもの、つまり心像のうちどのようなものを用いるべきかである。(92)
\\

\textbf{困難な状況に対してどのように挑むべきか}

困難な状況は人間の真価を示すものである。(135)

敵など近くに一人もいない。万事が平和に満ちている。見たまえ、私は殴られていないし、傷つけられていないし、だれかから逃げたりもしていない。(136)
\\

\textbf{人々に対して起こるべきでないこと、そして、人々の間では何が小事で、何が大事か}
人が虚偽を承認するときは、虚偽を承認することをのぞんでしたわけではない。

目の見えない人をかわいそうに思うように、善悪の判断ができないのを同情するべきだ。

正しい思考が奪われたとき、これが人間の滅亡だ。


\textbf{困難な状況に対してどのように挑むべきか}
すべての人に当てはまる原則がある。それは、物事を承認するのはそのとおりだと感じるからで、承認しないのはそうでないからだ。

泥棒や盗賊は、善と悪の判断を誤っているのだ。彼らに怒るのではなく、過失を教えてやるべきだ。

我々はどうして腹を立てる？奪い取られたものを大事に思うからだ。

不敗の者とは？それは意志とは関係ないどんなものにも惑わされない人だ。
